% ============================================================
% DRAFT — proposed edits to chapters/evaluation.tex, from the new
% context-budget sweep CSVs you added today, extended from depth 40
% out to depth 120 (ctx_1024_d120.csv, ctx_2048_d120.csv,
% ctx_8192_d120.csv; 3 runs per depth, every 4 depths from 4 to 120,
% 90 runs per context setting, 270 runs total).
%
% This supersedes the SHORTER depth-40 sweep that the current
% "Context window and the plan ceiling" paragraph, Table~\ref{tab:ctx}
% and the "Measuring the planning prompt..." paragraph are built on.
% This file is a draft for review, not a finished chapter section.
% Read it next to evaluation.tex, copy over whatever you're happy
% with, and edit the rest — nothing here has been written into the
% actual chapter file.
%
% IMPORTANT — scope check first, because this is NOT everything in
% sec:rq2_results: I went back and checked where every part of that
% section's data actually comes from. The opening paragraph,
% Table~\ref{tab:rq2_depth}, "Where execution failed," and "Probe
% chains" are all built from tour1_main.csv / tour2_main.csv, run on
% 4 October — approach/literal/mixed/probe chains at depths 1 to 12,
% local AND cloud. Today's new sweep (ctx_*_d120.csv, run 20 October)
% is a completely different, separate experiment: approach chains
% ONLY, local ONLY, depths 4 to 120, three context settings, and
% zero collisions, zero hallucinations, zero probes, outcome
% "completed" in all 270 rows. It does not touch, support, or
% contradict anything in the opening paragraph, Table~\ref{tab:rq2_depth},
% "Where execution failed," or "Probe chains" — those four pieces are
% a different test and do NOT need changing because of today's runs.
% If you have newer corpus data (a re-run of the depth-1-to-12
% approach/literal/mixed/probe test) sitting somewhere else, say so
% and I'll redo those four pieces too — but the three _d120 CSVs
% alone don't give me anything to change them with.
%
% What the new sweep DOES affect is exactly the part of
% sec:rq2_results that comes after "Probe chains": "Context window
% and the plan ceiling," Table~\ref{tab:ctx}, Figure~\ref{fig:ctx},
% and the "Measuring the planning prompt..." paragraph that follows
% the table. All four are rewritten below (E1-E4), and I've also
% carried the corrected RQ2 synthesis into the Discussion/Conclusion
% text (E5) for consistency.
%
% Blocks in this file, in the order they appear in evaluation.tex:
%   E1. Replace the "\paragraph{Context window and the plan ceiling.}"
%       paragraph (prose only).
%   E2. Replace Table~\ref{tab:ctx} with corrected, depth-120 figures.
%   E3. Figure~\ref{fig:ctx} needs regenerating from the three new
%       _d120 CSVs (it currently only shows data out to depth 40).
%       I haven't touched the image — see the note where this block
%       sits for what the regenerated plot should show.
%   E4. FULL rewrite of the "Measuring the planning prompt..."
%       paragraph. This is NOT a light edit any more — on checking
%       it against the corrected Table~\ref{tab:ctx} (E2), several of
%       its sentences are now simply wrong (see the note at E4 for
%       exactly which, and why). The \candidatetodo at the end stays,
%       strengthened: the new data makes the "was num_ctx even
%       enforced?" question more pressing, not less.
%   E5. Update to the \paragraph{RQ2.} paragraph — this already lives
%       in conclusion_move_discussion_from_evaluation.tex (the
%       Discussion section moved there), reproduced here only so this
%       file's E1-E5 numbering still makes sense on its own.
%
% Headline facts behind E1-E4, if you want to sanity-check any
% wording against the actual numbers:
%   - 0 hallucinations in all 270 runs, at every depth, every context
%     setting. outcome="completed" for all 270 — nothing halts, no
%     collision, no blocked step, no probe in this sweep.
%   - first_divergence always lands exactly where plan_steps stops:
%     the plan is a correct, in-order prefix right up to where it
%     stops, never a wrong substitution. This is truncation, not
%     hallucination.
%   - Every run is correct (visits_match=True) through depth 28, at
%     all three context settings.
%   - 1024: 3/3 correct at every depth 4-28, 2/3 at 32, back up to 3/3
%     at 36, then 0/63 correct from depth 40-120 (mean 24.6 steps,
%     26/63 empty plans).
%   - 2048: 3/3 correct at every depth 4-36 (no slip before 40), then
%     0/63 correct from depth 40-120 (mean 23.7 steps, 27/63 empty).
%   - 8192: 3/3 correct at every depth 4-32, 2/3 at 36, then 0/63
%     correct from depth 40-120 (mean 46.7 steps, 0/63 empty).
%   - So the three settings fail completely at the SAME depth (40),
%     despite an eightfold difference in context budget, and the
%     small dip-then-recover at depth 32/36 for 1024 and 8192 is
%     noise at n=3, not a trend.
% ============================================================


% ---------- E1. "Context window and the plan ceiling" paragraph ----------
% CURRENT (chapters/evaluation.tex, inside \mysection{Compositional depth
% and hallucination onset}, immediately after the "Probe chains."
% paragraph):
%   \paragraph{Context window and the plan ceiling.} Figure~\ref{fig:ctx}
%   shows the context sweep. With the configured \texttt{num_ctx} of 8192
%   the local planner produced correct plans for chains up to 36 instructions
%   in all three runs, and at 40 instructions it produced a 34-step plan in all three
%   runs, dropping the last six visits. The ceiling is therefore a plan-length
%   limit of the model, and not a context limit, since it appears at 8192 where
%   the plan uses only a fraction of the window.
% REPLACE WITH:

\paragraph{Context window and the plan ceiling.} I re-ran the context sweep and
extended it from depth 40 out to depth 120: three runs per depth, every four
depths from 4 to 120, 90 runs per context setting, 270 runs in total.
Figure~\ref{fig:ctx} shows the extended sweep. No hallucination showed up in any
of the 270 runs, at any depth or context setting. What actually happens is that
the planner stops generating further steps partway through an otherwise-correct,
in-order plan: \texttt{first\_divergence} always lands exactly on the step where
the plan stops, it never substitutes a wrong object, location or action. Every
context setting produced a correct plan in all three runs from depth 4 to 28.
There is some noise around depths 32--36 (one of three runs wrong for
\texttt{num\_ctx}=1024 at depth 32, and for \texttt{num\_ctx}=8192 at depth 36),
but by depth 40 every run at every context setting fails, and none of them
recover at any depth tested up to 120. So the ceiling is a plan-length limit on
the model's own generation, not a context-window limit: going from 1024 to 8192
tokens, an eightfold increase, does not move the depth where the plan starts to
fail. The context window does change one thing: how badly each run fails once it
starts. At depths 40 and above, \texttt{num\_ctx}=1024 and 2048 both produce a
completely empty plan in roughly 40\% of runs (26 of 63 and 27 of 63), while 8192
never produces an empty plan and keeps close to double the mean plan length of
the other two settings (46.7 steps against 23.7--24.6). A bigger context window
does not delay when the plan starts to degrade, but it does raise how much of a
degrading plan survives.


% ---------- E2. Table~\ref{tab:ctx} — corrected, extended to depth 120 ----------
% CURRENT:
%   \begin{table}[htbp]
%       \centering
%       \caption{Correct plans out of three runs, by context setting and chain depth (Gemma~4 E4B).}
%       \label{tab:ctx}
%       \begin{tabular}{lcccc}
%           \toprule
%           \texttt{num_ctx} & Depth 4--32 & 36 & 40 & Empty plans \\
%           \midrule
%           8192 & 3/3 at every depth & 3/3 & 0/3 (34 steps) & 0 \\
%           2048 & 3/3 at every depth & 2/3 & 0/3 & 1 (depth 40) \\
%           1024 & 3/3 at every depth & 2/3 & 1/3 & 1 (depth 40) \\
%           \bottomrule
%       \end{tabular}
%   \end{table}
% REPLACE WITH:

\begin{table}[htbp]
    \centering
    \caption{Correct plans out of three runs, by context setting and chain depth,
    extended to depth 120 (Gemma~4 E4B). The deep band (depth 40--120) summarises
    63 runs per context setting; no run in this band recovers to a correct plan at
    any depth.}
    \label{tab:ctx}
    \begin{tabular}{lccccc}
        \toprule
        \texttt{num\_ctx} & Depth 4--28 & 32 & 36 & Depth 40--120 (mean plan length) & Empty plans \\
        \midrule
        8192 & 3/3 at every depth & 3/3 & 2/3 & 0/63 correct, mean 46.7 steps & 0/63 \\
        2048 & 3/3 at every depth & 3/3 & 3/3 & 0/63 correct, mean 23.7 steps & 27/63 \\
        1024 & 3/3 at every depth & 2/3 & 3/3 & 0/63 correct, mean 24.6 steps & 26/63 \\
        \bottomrule
    \end{tabular}
\end{table}


% ---------- E3. Figure~\ref{fig:ctx} needs regenerating ----------
% CURRENT:
%   \begin{figure}[htbp]
%       \centering
%       \includegraphics[width=\textwidth]{figures/fig_ctx_sweep_v2.png}
%       \caption{(a) Plans with the correct visit order out of three runs, against chain depth, at three context settings (points are offset horizontally for legibility). (b) Projected tokens per planning call against depth, measured to depth 30, with the 1024 and 2048 limits marked.}
%       \label{fig:ctx}
%   \end{figure}
% The \includegraphics call and label can stay as-is; only the PNG itself
% (and, if you want, the caption) need updating, since fig_ctx_sweep_v2.png
% was built from the depth-40 sweep. Regenerate it from the three new
% _d120 CSVs so panel (a) runs out to depth 120 on the x-axis instead of
% stopping at 40 — that's the whole point of the new data, and the plot
% currently undersells it. Suggested caption update, if the regenerated
% plot keeps the same two-panel layout:
%
%   \caption{(a) Plans with the correct visit order out of three runs,
%   against chain depth from 4 to 120, at three context settings (points
%   are offset horizontally for legibility). Every setting fails fully
%   from depth 40 onward, with no recovery at any higher depth. (b)
%   Projected tokens per planning call against depth, measured to depth
%   30, with the 1024 and 2048 limits marked.}
%
% If you'd rather not touch the figure right now, at minimum change the
% caption so it doesn't claim the sweep stops at 40 while the table next
% to it reports results to 120 -- as it stands the table and figure would
% disagree on how far the data goes.


% ---------- E4. "Measuring the planning prompt..." paragraph — FULL rewrite ----------
% CURRENT (chapters/evaluation.tex, the paragraph right after Table~\ref{tab:ctx}):
%   Measuring the planning prompt gave 460 tokens at
%   depth 1 and 1\,782 at depth 30 (about 46 tokens per added instruction),
%    which is 22\% of the configured 8192-token window. For a mission that
%    includes replanning with three guidance turns, the largest prompt
%    measured was 1\,271 tokens at depth 30, and the largest revision prompt 1\,579 tokens.
%    Extrapolating the single-shot slope gives an overflow at about depth 170 for 8192 tokens and
%    about depth 2\,800 for the native 128K window, far beyond the 34-step ceiling.
%    The same projection predicts overflow at about depth 13 for 1024 and about depth
%    36 for 2048 tokens. Table~\ref{tab:ctx} does not follow this for 1024: plans stayed
%    correct up to depth 32 at a setting where the projection says the prompt should not fit.
%    The prediction at 2048 matches the observed drop at depth 36 only by coincidence,
%    since the same drop occurs at 8192. One empty plan at depth 40 appeared at each reduced
%    setting and none at 8192, which is suggestive but is one run each.
%    \candidatetodo{I could not verify that the server enforced the reduced context.
%    Run ollama ps during a sweep (or an overflow probe longer than 1024 tokens) and
%    state the loaded context length here before using the 1024 and 2048 results as evidence either way.}
%
% WHY this needs a full rewrite and not just a wording tweak: the last
% three sentences were written against the OLD, depth-40-only
% Table~\ref{tab:ctx}, and every one of them is now wrong against the
% corrected table in E2:
%   - "plans stayed correct up to depth 32" for 1024 — no: the corrected
%     table has 1024 at 2/3 (not 3/3) at depth 32.
%   - "the observed drop at depth 36 ... coincidence, since the same
%     drop occurs at 8192" for 2048 — no: the corrected table has 2048
%     at 3/3 (no drop at all) at depth 36. Only 8192 drops at 36.
%   - "one empty plan ... at each reduced setting ... one run each" —
%     no: the deep band (depth 40-120) is now 63 runs per setting, with
%     26/63 and 27/63 empty for 1024/2048, not one run total.
% REPLACE THE ENTIRE PARAGRAPH WITH:

Measuring the planning prompt gave 460 tokens at depth 1 and 1\,782 at depth 30
(about 46 tokens per added instruction), which is 22\% of the configured
8192-token window. For a mission that includes replanning with three guidance
turns, the largest prompt measured was 1\,271 tokens at depth 30, and the
largest revision prompt 1\,579 tokens. If I extrapolate that slope, the prompt
should run out of room at around depth 13 for \texttt{num\_ctx}=1024, depth 36
for 2048, and depth 170 for 8192 (about depth 2\,800 for the native 128K
window). None of that matches what the depth-120 sweep actually shows. Every
context setting holds correct plans through at least depth 28 and then fails
completely from depth 40 onward, no matter what \texttt{num\_ctx} is set to.
The 1024 setting stays correct more than twice as deep as its own token budget
says it should, while 8192 fails more than four times earlier than its
projection predicts. If the model were genuinely running out of context room,
I would expect it to fail close to its own projected depth, and to fail later
as the window grows. Instead all three settings fail at almost exactly the
same depth. Together with the identical depth-40 failure point, this tells me
the ceiling is about how long a plan the model will generate, not how much
context it has to read.
\candidatetodo{I could not verify that the server actually enforced the
reduced context size at all, and this result makes that worth checking more,
not less: if \texttt{num\_ctx}=1024 and 8192 fail at the same depth despite an
eightfold difference in token budget, that is consistent with both "context
size genuinely does not matter here" and "the server never applied the
smaller setting in the first place." Run \texttt{ollama ps} during a sweep
(or an overflow probe longer than 1024 tokens) and state the loaded context
length here before using the 1024 and 2048 results as evidence of anything
about context size specifically.}


% ---------- E5. Discussion/Conclusion \paragraph{RQ2.} ----------
% This paragraph already lives in
% conclusion_move_discussion_from_evaluation.tex (block M2), since you
% asked for the whole Discussion section to move into conclusion.tex.
% It is unchanged by this revision — reproduced here only so this
% file's numbering is self-contained; don't apply it twice.
%
%   \paragraph{RQ2.} In the virtual tests no hallucinated step appeared at any
%   depth tested, including the extended context sweep run out to 120 chained
%   instructions: across all 270 of those runs, the planner never substitutes a
%   wrong object, location or action. Compositional depth does have a hard
%   limit, but the failure mode is truncation, not hallucination: from depth 40
%   onward, every run at every context setting tested (1024, 2048 and 8192
%   tokens) stops generating steps partway through an otherwise-correct plan,
%   and none recovers at any depth up to 120. Context window size does not move
%   this onset, but it does change its severity: the smaller context settings
%   produce a completely empty plan in roughly 40\% of deep runs, while the
%   largest setting never does and sustains close to double the plan length on
%   average, which points to a generation-length ceiling in the planner's own
%   decoding behaviour rather than the model running out of context-window
%   capacity to track the room and the command. Neither this truncation nor the
%   silent dropping of unresolvable steps described above is caught by the
%   pipeline at plan time, which argues for a capability and coverage check on
%   the plan before departure as the most useful next improvement.


% ---------- On moving "Reference comparison against cloud models" ----------
% You floated moving sec:cloud_comparison into the conclusion too, then
% talked yourself back out of it mid-message. My take: leave it in
% evaluation.tex. It's mostly results (the 63-vs-105 correctness match,
% the latency numbers, the probe-dropping behaviour matching locally),
% which belongs with the other evidence, not in the Conclusion chapter.
% The one genuinely interpretive bit at its end — "no cloud run was made
% on the physical rover" and the depth-40 sweep never being run on the
% cloud model — is a limitation, and I already folded that into
% conclusion_move_discussion_from_evaluation.tex as its own "Reference
% comparison scope" paragraph under Limitations. So the synthesis half
% of this section is already where you wanted it; the data half should
% stay put.
